%!TeX encoding = UTF-8
\documentclass{article}
\usepackage[pdftex]{graphicx} %for embedding images
\usepackage{url} %for proper url entries
\usepackage[bookmarks, colorlinks=false, pdfborder={0 0 0}, pdftitle={Final Report - Recognition Spring 2021}, pdfauthor={Nhut-Nam Le}, pdfsubject={Recognition}, pdfkeywords={report, exercises}]{hyperref} %for creating links in the pdf version and other additional pdf attributes, no effect on the printed document
%\usepackage[final]{pdfpages} %for embedding another pdf, remove if not required
\usepackage[utf8]{inputenc}
\usepackage[vietnamese]{babel}
\usepackage{float}
\usepackage{fancyhdr}
\usepackage{pythonhighlight}
\usepackage[left=3cm, right=3cm, top=2cm, bottom=2cm]{geometry}
\usepackage{parskip}
\usepackage{tikz}
\usepackage{hyperref}
\usepackage[]{algorithm2e}
\usepackage[noend]{algpseudocode}
\usepackage{amsmath}
\usepackage{amsfonts}

\usepackage{listings}
\usepackage{color}

\definecolor{dkgreen}{rgb}{0,0.6,0}
\definecolor{gray}{rgb}{0.5,0.5,0.5}
\definecolor{mauve}{rgb}{0.58,0,0.82}

\newcommand\T{\rule{0pt}{2.6ex}}       % Top strut
\newcommand\B{\rule[-1.2ex]{0pt}{0pt}} % Bottom strut

\lstset{frame=tb,
	language=Java,
	aboveskip=3mm,
	belowskip=3mm,
	showstringspaces=false,
	columns=flexible,
	basicstyle={\small\ttfamily},
	numbers=none,
	numberstyle=\tiny\color{gray},
	keywordstyle=\color{blue},
	commentstyle=\color{dkgreen},
	stringstyle=\color{mauve},
	breaklines=true,
	breakatwhitespace=true,
	tabsize=3
}

\setlength{\parindent}{15pt}
\setlength{\headheight}{15.2pt}
\pagestyle{fancy}
\lhead[<even output>]{NHẬN DẠNG}
\rhead[<even output>]{ĐỀ CƯƠNG NGHIÊN CỨU}
\title{research-outline}
\author{Nhut-Nam Le}
\date{2021}

\usepackage{enumitem}
\setlistdepth{20}
\renewlist{itemize}{itemize}{20}

% initially, use dots for all levels
\setlist[itemize]{label=$\cdot$}

% customize the first 3 levels
\setlist[itemize,1]{label=\textbullet}
\setlist[itemize,2]{label=--}
\setlist[itemize,3]{label=*}
\begin{document}
	\begin{titlepage}
		\begin{center}
			% Top of the page
			\large{\textbf{ĐẠI HỌC KHOA HỌC TỰ NHIÊN, ĐHQG-HCM\\KHOA CÔNG NGHỆ THÔNG TIN\\BỘ MÔN KHOA HỌC MÁY TÍNH}}\\
			\includegraphics[width=0.75\textwidth]{images/khtn.png}\\
			% Title
			\huge \textbf{NHẬN DẠNG}\\[0.1in]
			\huge \textbf{ĐỀ CƯƠNG NGHIÊN CỨU}\\[0.1in]
			\vfill
			\normalsize
			% Submitted by
			\normalsize
			% Lecturers
			\textbf{Giảng viên lý thuyết}\\
			{\textbf{PGS. TS.} Lê Hoàng Thái}\\[0.1in]
			% Teacher Assistant
			\textbf{Giảng viên hướng dẫn}\\
			\vspace{0.1in}
			{Lê Thanh Phong, Nguyễn Ngọc Thảo}\\[0.1in]
			\textbf{Sinh viên thực hiện} \\
			\vspace{0.1in}
			{Nguyễn Hoàng Đức, Lê Nhựt Nam, Nguyễn Viết Dũng}\\[0.1in]
			% Date time when written report
			\vfill
			Tháng 4 năm 2021
		\end{center}
	\end{titlepage}
	\newpage
	% End Title4
	
	\pagenumbering{roman} %numbering before main content starts
	\cleardoublepage
	%\pagebreak
	
	\newpage
	\tableofcontents
	\newpage
	\pagenumbering{arabic} %reset numbering to normal for the main content
	\setcounter{secnumdepth}{0}
	
	\section{1. THÔNG TIN NHÓM}
	\begin{itemize}
		\item Thành viên 01: Nguyễn Hoàng Đức - 18120018
		\item Thành viên 02: Lê Nhựt Nam - 18120061
		\item Thành viên 03: Nguyễn Viết Dũng - 18120167
	\end{itemize}
	
	\section{2. THÔNG TIN ĐỀ TÀI}
	\subsection{Tên đề tài}
	\begin{itemize}
		\item Tên đề tài (Tiếng Việt): Sinh trắc học giọng nói
		\item Tên đề tài (Tiếng Anh): Voice Biometrics
	\end{itemize}
	\subsection{Nguồn tham khảo}
	\begin{itemize}
		\item Sách: Chương 8: Voice Biometrics, Handbook of Biometric, Anil K. Jain, Patrick Flynn, Arun A.
		Ross
		\item Paper with code: Speaker Recognition
	\end{itemize}
	Danh sách các bài báo khoa học tham khảo:
	\begin{itemize}
		\item Deep neural networks for small footprint text-dependent speaker
		verification, Ehsan Variani (Johns Hopkins Univ., Baltimore, MD, USA), Xin Lei (Google Inc.,
		USA), Erik McDermott (Google Inc., USA), Ignacio Lopez Moreno (Google Inc, Mountain View,
		CA, US), Javier Gonzalez-Dominguez (Google Inc., USA), 2014 IEEE International Conference on Acoustics, Speech and Signal
		Processing (ICASSP), Florence, Italy
		\item Multi-Task Learning for Text-Dependent Speaker Verification, Nanxin Chen, Yanmin Qian, Kai Yu (Shanghai Jiao Tong University, China), INTERSPEECH 2015 16th Annual Conference of the International Speech Communication Association, Dresden, Germany, 2015
		\item X-Vectors: Robust DNN Embeddings for Speaker Recognition, David Snyder, Daniel Garcia-Romero, Gregory Sell, Daniel Povey, Sanjeev Khudanpur (Center for Language and Speech Processing \& Human Language Technology Center of Excellence, The Johns Hopkins University, Baltimore, MD, USA), 2018 IEEE International Conference on Acoustics, Speech and Signal
		Processing (ICASSP), Calgary, AB, Canada, 2018
		\item Mirco Ravanelli and Yoshua Bengio. Speaker recognition from raw waveform with sincnet, 2019.
	\end{itemize}
	\subsection{Từ khóa}
	\begin{itemize}
		\item Tên từ khóa (Tiếng Việt): Nhận dạng sinh trắc học, Nhận dạng giọng nói, Nhận dạng người nói, Mạng Neural Tích chập, Mẫu thô, Xác minh người nói, Định danh người nói
		\item 	Tên từ khóa (Tiếng Anh): Biometric Recognition, Voice Recognition, Speaker Recognition, Convolutional Neural Networks, Raw Samples, Speaker Verification, Speaker Identification
	\end{itemize}
	
	\subsection{Nội dung trình bày}
	\qquad Nội dung trình bày trước lớp gồm 2 phần: trình bày nội dung tìm hiểu được từ chương 08 - Voice Biometrics sách Handbook of Biometric, trình bày những phương pháp state-of-the-art về lĩnh vực nhận dạng giọng nói trong thời gian gần đây.
	
	Phần trình bày nội dung tìm hiểu được từ sách
	\begin{itemize}
		\item Giới thiệu chung
		\item Những thông tin nhận dạng trong tín hiệu giọng nói
		\item Rút trích đặc trưng và phân tách
		\item Hai công nghệ chính của lĩnh vực nhận dạng giọng nói
	\end{itemize}
	Phần trình bày những phương pháp state-of-the-art về lĩnh vực nhận dạng giọng nói
	\begin{itemize}
		\item Giới thiệu
		\item Động lực nghiên cứu khoa học
		\item Phát biểu bài toán
		\item Các công trình tiêu biểu (khoảng 3 công trình)
		\begin{itemize}
			\item Deep neural networks for small footprint text-dependent speaker verification - Đại diện cho d-vectors (Deep Vectors)
			\item Multi-Task Learning for Text-Dependent Speaker Verification - Đại diện cho j-vectors
			\item X-Vectors: Robust DNN Embeddings for Speaker Recognition - Đại diện cho x-vectors
			\item Speaker recognition from raw waveform with sincnet - Đại diện cho Phân lớp người nói
		\end{itemize}
		\item So sánh d-vectors, j-vectors và x-vectors
		\item Demo
		\item Tài liệu tham khảo
	\end{itemize}
	\subsection{Nội dung báo cáo}
	\begin{itemize}
		\item Thông tin nhóm
		\item Thông tin đề tài
		\item Nội dung phân công
		\item Nội dung báo cáo
		\begin{itemize}
			\item Nội dung tìm hiểu từ sách Handbook of Biometric
			\begin{itemize}
				\item Giới thiệu chung
				\item Thông tin nhận dạng trong tín hiệu giọng nói
				\item Rút trích đặc trưng và phân tách thông tin
				\begin{itemize}
					\item Phân tích theo từng đoạn ngắn
					\item Tham số hóa
					\item Phân tách ngữ âm và tách từ
					\item Phân tách ngữ điệu
				\end{itemize}
				\item Công nghệ nhận dạng giọng nói phụ thuộc văn bản
				\item Công nghệ nhận dạng giọng nói không phụ thuộc văn bản
			\end{itemize}
			\item Các phương pháp SOTA
			\begin{itemize}
				\item Giới thiệu chung định hướng gần đây
				\item Động lực nghiên cứu
				\item Phát biểu bài toán: Đầu vào (Input), Đầu ra (Output)
				\item Kho ngữ liệu
				\item Kiểm định mô hình nhận dạng giọng nói
				\item Các độ đo thường dùng trong nhận dạng giọng nói
				\item Các công trình tiêu biểu
				\begin{itemize}
					\item Trong tác vụ rút trích đặc trưng
					\begin{itemize}
						\item d-vectors
						\item j-vectors
						\item x-vectors
						\item So sánh d-vectors, j-vectors và x-vectors
					\end{itemize}
					\item Trong phân lớp người nói
					\begin{itemize}
						\item Variational autoencoder
						\item Multi-domain features
						\item SincNet
					\end{itemize}
				\end{itemize}
			\end{itemize}
			\item Thực nghiệm: Sử dụng SincNet trong xác minh và định danh người nói
			\begin{itemize}
				\item Chuẩn bị dữ liệu
				\item Xây dựng mô hình
				\item Đăng ký
				\item Đánh giá
			\end{itemize}
		\end{itemize}
	\end{itemize}
	\subsection{Xây dựng demo cho chủ đề nghiên cứu}
	\begin{itemize}
		\item Phương pháp giải quyết vấn đề
		\item Dữ liệu thực nghiệm
		\begin{itemize}
			\item Tên tập dữ liệu: TIMIT Acoustic-Phonetic Continuous Speech Corpus
			\item Tác giả: 	John S. Garofolo, Lori F. Lamel, William M. Fisher, Jonathan G. Fiscus, David S. Pallett, Nancy L. Dahlgren, Victor Zue
			\item Năm: 1993
			\item DCMI Type(s): âm thanh - sound
			\item Sample Type: 1-channel pcm (mono channel)
			\item Nguồn thu dữ liệu: microphone speech
			\item Ngôn ngữ: tiếng Anh
			\item Kích thước: 630 người nói từ 8 đặc trưng giọng Anh Mỹ, 1,3GB
		\end{itemize}
		\item Thực nghiệm và đánh giá
	\end{itemize}
	
	\section{3.	THÔNG TIN TỰ ĐÁNH GIÁ TIẾN ĐỘ}
	\subsection{Những công việc đã thực hiện được}
	
	\subsection{Những công việc chưa thực hiện được}
\end{document}